
\chapter{Problem Formulation and Proposed Solution}

\section{Figures and Tables}

Figures are conveniently included using postscript format.  If you are
generating a figure in a software, please check if the software
supports writing to a postscript or a PDF format. This format is loss
less vector format and with reproduce in any magnification without any
pixelation. Make sure to write it to an ``Encapsulated Post-script''or
.eps format.


\begin{figure}[!htbp]
  \centering
   \includegraphics[width=0.65\textwidth]{Fig1}
    \caption[UML Class diagram]{UML Class diagram showing an example of the Facade design pattern. The caption of the figure goes here. A shorter caption can be written in square brackets to identify it in the list of figures if required.}
    \label{fig:1} 
\end{figure}

Figures should be given a label and which can be used to refer to them
in the running text using \verb|\ref{}| command. Figure~\ref{fig:1}
describes the process flow sheet of the Facade design pattern used in
this study. 
Two simple graphs are shown in Figure~\ref{fig:2}. 
Multiple figures with same caption can be arranged as shown in Figure~\ref{fig:2} and they are referred in text such as Figure~\ref{fig:2a} and Figure~\ref{fig:2b}. The graphs should be drawn at appropriate places with center alignment and it should be referred in text.

\begin{figure}[!htbp]
     \centering
     \begin{subfigure}{0.45\textwidth}
         \centering
         \includegraphics[width=\textwidth]{Fig2a}
         \caption{Figure2a}
         \label{fig:2a}
     \end{subfigure}
     \hfill
     \begin{subfigure}{0.45\textwidth}
         \centering
         \includegraphics[width=\textwidth]{Fig2b}
         \caption{Figure2b}
         \label{fig:2b}
     \end{subfigure}
        \caption{Two simple graphs}
        \label{fig:2}
\end{figure}


The format for tables is given below. The table must referred in
the text as Table~\ref{tab:1}. The title of the table with table number should be written at the top of the table with center aligned as shown below in Table~\ref{tab:1}.

\begin{table}[!htbp]
    \centering
    \caption{A table shows an example of multirow and multicolumn.}
    \label{tab:1}
\begin{tabular}{|l|l|l|l|}\hline
  \multirow{10}{*}{numeric literals} & \multirow{5}{*}{integers} & in decimal & \verb|8743| \\ \cline{3-4}
  & & \multirow{1}{*}{in octal} & \verb|0o7464| \\ \cline{2-4}
  & & \multirow{2}{*}{in hexadecimal} & \verb|0x5A0FF| \\ \cline{4-4}
  & & & \verb|0xE0F2| \\ \cline{2-4}
  & \multirow{4}{*}{fractionals} & \multirow{4}{*}{in decimal} & \verb|140.58| \\ \cline{4-4}
  & & & \verb|8.04e7| \\ \cline{4-4}
  & & & \verb|0.347E+12| \\ \cline{4-4}
  & & & \verb|47e22| \\ \cline{1-4}
  \multicolumn{3}{|l|}{\multirow{3}{*}{char literals}} & \verb|`H'| \\ \cline{4-4}
  \multicolumn{3}{|l|}{} & \verb|`\n'| \\ \cline{4-4}          %% here
  \multicolumn{3}{|l|}{} & \verb|`\x65'| \\ \cline{1-4}        %% here
  \multicolumn{3}{|l|}{\multirow{2}{*}{string literals}} & \verb|``bom dia"| \\ \cline{4-4}
  \multicolumn{3}{|l|}{} & \verb|``ouro preto\nmg"| \\ \cline{1-4}          %% here
\end{tabular}
\end{table}

A simple table is shown in Table~\ref{tab:2}. One may use the online latex table generator to get the latex code for required table: \url{https://www.tablesgenerator.com/} 

\begin{table}[!htbp]
\centering
    \caption{Simple table.}
    \label{tab:2}
    \begin{tabular}{|l|c|r|} % <-- Alignments: 1st column left, 2nd middle and 3rd right, with vertical lines in between
    \hline
      \textbf{Value 1} & \textbf{Value 2} & \textbf{Value 3}\\
      $\alpha$ & $\beta$ & $\gamma$ \\
      \hline
      1 & 1110.1 & a\\ \hline
      2 & 10.1 & b\\ \hline
      3 & 23.113231 & c\\ \hline
    \end{tabular}
\end{table}


\section{Code}

A simple JAVA code is given below:

\lstset{frame=tb,language=Java} %%% may change the language according to the code
\begin{lstlisting}
// Hello.java
import javax.swing.JApplet;
import java.awt.Graphics;

public class Hello extends JApplet {
    public void paintComponent(Graphics g) {
        g.drawString("Hello, world!", 65, 95);
    }    
}
\end{lstlisting}




%%% Local Variables: 
%%% mode: latex
%%% TeX-master: "../mainrep"
%%% End: 
